\chapter{Desenvolvimento}\label{cap:desenvolvimento}
\addcontentsline{toc}{chapter}{Desenvolvimento}

Este capítulo descreve a implementação efetivamente utilizada para executar o protocolo definido no Capítulo~\ref{cap:metodologia}. O artefato resultante é um repositório controlador que adquire e verifica o corpus, prepara entradas sanitizadas, executa as seis condições experimentais, valida as respostas, preserva todas as tentativas e produz evidências auditáveis. Diferentemente de uma especificação anterior à coleta, o texto a seguir registra as versões, os controles, os formatos e os desvios observados durante a execução.

A coleta das seis condições foi concluída sobre os 26 repositórios do RealVuln v1.0. Foram selecionadas 364 unidades válidas: 52 execuções SAST, correspondentes a uma execução de C1 e C2 por alvo, e 312 execuções de IA, correspondentes a três repetições de C3--C6 por alvo. As retentativas não foram sobrescritas; somadas às unidades selecionadas, produziram 433 manifestos ou tentativas preservadas. Esses números caracterizam a completude operacional do artefato. As medidas de detecção e a comparação com o \textit{ground truth} pertencem ao Capítulo~\ref{cap:avaliacao}.

\section{Visão Geral da Implementação}\label{sec:visao-geral-desenvolvimento}

O controlador foi implementado predominantemente em Python e PowerShell. Python concentra os modelos de dados, a validação dos contratos, a preparação do corpus, a fila, a execução e a normalização de C1 e C2. PowerShell coordena a aquisição, o Docker Desktop e as interfaces de linha de comando do Cursor e do Codex. Os projetos vulneráveis foram tratados como dados: não foram importados, testados nem executados pelo controlador.

O fluxo foi dividido em três domínios. O domínio de entrada contém as árvores sanitizadas e os \textit{payloads} de código entregues às ferramentas. O domínio de controle contém configurações congeladas, filas, \textit{prompts}, esquemas, adaptadores e áreas de trabalho. O domínio de avaliação contém o \textit{ground truth} e o validador oficial do RealVuln. Durante a coleta, os executores de C1--C6 não consultaram o domínio de avaliação; as auditorias das condições de IA registraram explicitamente \texttt{oracle\_consultado=false}. Essa separação materializa o cegamento descrito na Seção~\ref{sec:repositorio-experimental}.

A cadeia implementada pode ser resumida em: aquisição das revisões fixadas; sanitização e inventário; congelamento das configurações; execução serial e preservação das tentativas; validação sintática e estrutural; auditoria de fechamento; e, em etapa separada, pontuação contra o \textit{oracle}. A separação entre coleta e pontuação impede que rótulos do benchmark influenciem a geração ou a normalização dos achados.

\section{Organização do Repositório Controlador}\label{sec:organizacao-repositorio-desenvolvimento}

A Tabela~\ref{tab:estrutura-repositorio-experimental} apresenta a estrutura efetivamente adotada. Três diretórios possuem papéis distintos na cadeia de custódia: \texttt{execucoes/} mantém o estado e as áreas mínimas de trabalho; \texttt{resultados/} contém os registros primários por tentativa; e \texttt{evidencias/} contém sínteses derivadas e inventários ligados aos registros primários por SHA-256.

\begin{table}[htb]
\ABNTEXfontereduzida
\caption{Estrutura efetiva do repositório controlador}
\label{tab:estrutura-repositorio-experimental}
\centering
\begin{tabularx}{\textwidth}{|l|X|X|}
\hline
\textbf{Diretório} & \textbf{Conteúdo} & \textbf{Regra de acesso e uso} \\
\hline
\texttt{config/} & \textit{Locks} de fontes, corpus, fila, sanitização, agentes, \textit{prompt} e esquema de saída & Contratos congelados; alterações posteriores invalidam a comparação com a coleta realizada \\
\hline
\texttt{benchmark/} & \textit{Cache} local reconstruível do RealVuln e dos repositórios de origem & Usado somente na aquisição; não é entregue às ferramentas nem redistribuído \\
\hline
\texttt{alvos/} & Árvores sanitizadas identificadas como \texttt{ALVO-NNNN} & Tratadas como dados e verificadas por inventário antes do uso \\
\hline
\texttt{oracle/} & Manifesto, \textit{ground truth} e validador oficiais & Segregado da coleta; acesso reservado à validação do benchmark e à pontuação posterior \\
\hline
\texttt{runner/} e \texttt{scripts/} & Núcleo Python, adaptadores e orquestradores PowerShell & Executam e validam o protocolo sem completar campos com informação do \textit{oracle} \\
\hline
\texttt{docker/} & Imagem, dependências fixadas e cópia local das regras Semgrep & Usado por C1 e C2; o pacote de regras é reconstruível, mas não é redistribuído \\
\hline
\texttt{execucoes/} & Fila, alertas híbridos congelados e áreas de trabalho das chamadas & Estado mutável do controlador; não é fonte única de métricas \\
\hline
\texttt{resultados/} & Saídas brutas, normalizadas, eventos, respostas e manifestos por tentativa & Registro primário e imutável; nenhuma tentativa é reutilizada ou sobrescrita \\
\hline
\texttt{evidencias/} & Inventários, auditorias de fechamento e sínteses derivadas & Evidência publicada, vinculada aos primários por hashes \\
\hline
\end{tabularx}
\legend{Fonte: Elaborado pelo autor (2026).}
\end{table}

Em C1 e C2, cada tentativa foi gravada em \path{resultados/CN-ALVO-NNNN-R01/tentativa-NNN/}. Em C3--C6, o caminho equivalente foi dividido por finalidade, por exemplo \path{resultados/ia/coleta/CN-ALVO-NNNN-RNN/tentativa-NNN/}. A finalidade \texttt{piloto} permaneceu em um \textit{namespace} separado e, por isso, não pôde ocupar uma unidade da coleta principal.

A publicação do artefato considerou também as licenças e o conteúdo sensível dos projetos de terceiros. Dezesseis árvores sanitizadas, com licença permissiva ou copyleft identificada, foram incluídas integralmente em \texttt{alvos/corpus-v1/}. Os outros dez alvos permaneceram reconstruíveis por URL, \textit{commit}, tamanho e SHA-256 de cada arquivo. Nove não apresentavam licença explícita na revisão adquirida e um continha uma chave privada de demonstração. O inventário publicado cobre ainda \texttt{oracle/}, \texttt{execucoes/} e \texttt{resultados/}, permitindo verificar a integridade do conjunto redistribuído sem publicar material não autorizado.

\section{Preparação do Corpus Experimental}\label{sec:preparacao-corpus}

\subsection{Fixação do RealVuln e verificação de integridade}\label{subsec:fixacao-realvuln}

O RealVuln foi fixado na tag \texttt{v1.0}, cujo objeto anotado é identificado pelo prefixo \texttt{aa9f7321c8c5} e cujo \textit{commit} resolvido possui o prefixo \texttt{d98e9fc91273}. O arquivo \texttt{config/corpus-realvuln-v1.lock.json} registra a URL e o \textit{commit} completo dos 26 projetos. Três URLs de origem estavam indisponíveis durante a aquisição; nesses casos foram usados espelhos previamente autorizados apenas como transporte, com aceitação condicionada à presença do mesmo objeto Git fixado no manifesto.

O script de aquisição obteve cada revisão sem executar \textit{hooks}, submódulos, testes ou arquivos do alvo. Antes da promoção das áreas temporárias, foram conferidos o manifesto, os \textit{commits}, os arquivos regulares permitidos e os hashes da árvore produzida. O validador oficial \texttt{validate\_gt.py} terminou com código zero sobre os 26 arquivos de \textit{ground truth}, totalizando 817 entradas: 697 vulnerabilidades e 120 armadilhas de falso positivo. O artefato preserva tanto essa saída literal quanto uma auditoria por repositório.

\subsection{Sanitização e identificadores opacos}\label{subsec:sanitizacao-alvos}

A política \texttt{sanitizacao-v1} removeu metadados \texttt{.git}, ambientes virtuais, dependências instaladas, \textit{caches}, artefatos de compilação, \textit{ground truth}, \textit{oracle}, soluções, \textit{walkthroughs} e \textit{write-ups}. Apenas objetos Git com modo de arquivo regular foram aceitos; três ligações simbólicas conhecidas foram omitidas sem que seus destinos fossem seguidos. A política foi congelada em \texttt{config/politica-sanitizacao-v1.json}.

Cada projeto recebeu um identificador opaco entre \texttt{ALVO-0001} e \texttt{ALVO-0026}. A correspondência entre esses identificadores e os projetos originais permaneceu no \textit{lock} do corpus, fora das áreas de trabalho dos agentes. Para cada alvo, o preparador produziu uma lista ordenada contendo caminho relativo, tamanho e SHA-256 de cada arquivo, além de um SHA-256 canônico da árvore. O executor verificou esse identificador antes e depois de C1 e C2; os executores de IA reconstruíram o \textit{payload} apenas com arquivos Python presentes no inventário congelado.

Essa estratégia também evitou depender de data de modificação ou de outros metadados do sistema de arquivos. A igualdade passou a ser definida pelos bytes e pelos caminhos relativos. As evidências em \texttt{evidencias/corpus-realvuln-v1/} permitem reconstruir e comparar todos os 26 alvos, inclusive aqueles que não puderam ser redistribuídos integralmente.

\section{Arquitetura do \textit{Harness}}\label{sec:arquitetura-harness}

\subsection{Módulos e responsabilidades}\label{subsec:modulos-harness}

A Tabela~\ref{tab:modulos-harness} relaciona os módulos implementados às suas responsabilidades. A decomposição separa aquisição, execução e avaliação, de modo que o código que entrega entradas às ferramentas não precise consultar o \textit{ground truth}.

\begin{table}[htb]
\ABNTEXfontereduzida
\caption{Módulos implementados no \textit{harness}}
\label{tab:modulos-harness}
\centering
\begin{tabularx}{\textwidth}{|l|X|X|}
\hline
\textbf{Componente} & \textbf{Responsabilidade implementada} & \textbf{Saída principal} \\
\hline
\texttt{aquisicao.py} e \texttt{corpus.py} & Validar fontes, revisões, políticas, espelhos e \textit{ground truth} & \textit{Locks}, auditorias e manifesto do corpus \\
\hline
\texttt{preparacao.py} & Copiar apenas arquivos permitidos e calcular inventários canônicos & Árvore sanitizada e inventário SHA-256 \\
\hline
\texttt{fila.py} e \texttt{coleta.py} & Controlar estados, exclusão mútua, retomada e fechamento de C1/C2 & Fila e histórico por tentativa \\
\hline
\texttt{executor\_sast.py} & Executar o scanner, medir duração, preservar bytes e fechar manifestos & Bruto, normalizado e manifesto terminal \\
\hline
\texttt{adaptadores/} & Converter Bandit e Semgrep para o esquema comum sem consultar o \textit{oracle} & Lista de achados normalizados \\
\hline
\texttt{scripts/ia/} & Preparar \textit{payloads}, chamar Cursor ou Codex e validar eventos e respostas & Resposta, eventos, metadados e manifesto de IA \\
\hline
\texttt{scripts/*-lote.ps1} & Planejar, executar serialmente, pular unidades válidas e retomar falhas & Resumo de lote e novas tentativas \\
\hline
\end{tabularx}
\legend{Fonte: Elaborado pelo autor (2026).}
\end{table}

Cada unidade recebeu um identificador formado por condição, alvo e repetição, como \texttt{C4-ALVO-0001-R02}. Em C1 e C2, a fila persistente admite os estados \texttt{pendente}, \texttt{em\_execucao}, \texttt{concluida} e \texttt{falha}, com histórico de transições e \textit{lock} de exclusão mútua. Em C3--C6, o lote revalida a tentativa existente antes de considerá-la concluída. Uma falha cria a próxima pasta \texttt{tentativa-NNN}; a pasta anterior permanece intacta.

Os registros terminais somente foram promovidos depois da gravação dos artefatos obrigatórios. Arquivos em produção receberam o sufixo \texttt{.part} e foram renomeados após o encerramento, enquanto JSONs de controle foram gravados de forma atômica. Os manifestos incluem hashes dos artefatos, horários, duração monotônica ou duração informada pelo adaptador, código de saída, versão observada, entrada e configuração. Esse desenho permitiu retomar as coletas sem confundir falha técnica com relatório vazio.

\subsection{Esquemas e preservação dos dados}\label{subsec:esquema-comum-achados}

As saídas nativas de Bandit e Semgrep foram preservadas primeiro em \texttt{bruto.json}. Somente depois o adaptador gerou \texttt{normalizado.json}. A Tabela~\ref{tab:esquema-achado-comum} resume o esquema dos achados SAST. Valores originais e normalizados coexistem; quando não existe mapeamento fechado, o campo normalizado permanece nulo.

\begin{table}[htb]
\ABNTEXfontereduzida
\caption{Esquema comum dos achados SAST normalizados}
\label{tab:esquema-achado-comum}
\centering
\begin{tabularx}{\textwidth}{|l|X|X|}
\hline
\textbf{Grupo} & \textbf{Campos principais} & \textbf{Tratamento} \\
\hline
Proveniência & condição, ferramenta, alvo, repetição, execução e hash da saída bruta & Obrigatórios e vinculados ao manifesto \\
\hline
Localização & caminho relativo, linha inicial e linha final & Normalização sintática; ausências não são inferidas \\
\hline
Classificação & regra, CWE, severidade e confiança, nas formas original e normalizada & Conversão apenas por mapas fechados \\
\hline
Relatório & descrição, evidência, recomendação e texto original & Conteúdo preservado sem complementação pelo controlador \\
\hline
Operação & início, término, duração, código de saída, estado e artefatos & Registrado no manifesto da tentativa \\
\hline
\end{tabularx}
\legend{Fonte: Elaborado pelo autor (2026).}
\end{table}

As respostas de C3--C6 utilizaram um contrato próprio, comum aos dois produtos. A raiz deve ser um \textit{array} JSON e cada item contém exatamente \texttt{arquivo}, \texttt{linha\_inicial}, \texttt{linha\_final}, \texttt{cwe}, \texttt{severidade}, \texttt{descricao} e \texttt{recomendacao}. Propriedades adicionais, caminhos ausentes do \textit{payload}, intervalos fora do arquivo, CWE malformada e severidade fora da enumeração invalidam a tentativa. A resposta bruta, contudo, permanece preservada mesmo quando a validação falha.

\section{Execução das Condições SAST (C1 e C2)}\label{sec:execucao-sast-desenvolvimento}

\subsection{Imagem Docker e isolamento}\label{subsec:construcao-imagem-docker}

C1 e C2 foram executadas com uma imagem experimental para \texttt{linux/amd64}, baseada em Python 3.12.13 sobre Debian Bookworm e fixada por \textit{digest}. A imagem continha Bandit 1.9.4 e Semgrep 1.172.0. As 69 dependências Python foram instaladas exclusivamente a partir de \textit{wheels} com versão e SHA-256 registrados. O corpus, os resultados e o \textit{oracle} não foram copiados para a imagem.

Cada tentativa utilizou um contêiner exclusivo com rede desativada, sistema de arquivos raiz somente leitura, usuário \texttt{10001:10001}, todas as \textit{capabilities} removidas, \texttt{no-new-privileges}, 2 CPUs, 3~GB de memória, 256 processos e \textit{tmpfs} de 64~MB com \texttt{noexec}, \texttt{nosuid} e \texttt{nodev}. A entrada foi montada como somente leitura e somente a pasta de saída recebeu escrita. O \texttt{docker.sock} não foi montado. O limite de cada scanner foi fixado em 900 segundos e as execuções ocorreram serialmente.

\subsection{Adaptador do Bandit}\label{subsec:execucao-bandit}

O Bandit 1.9.4 foi invocado recursivamente sobre \texttt{/entrada}, sem filtro posterior de severidade ou confiança:

\begin{verbatim}
bandit --recursive /entrada --format json \
  --output /saida/bruto.json.part
\end{verbatim}

Os códigos zero e um foram aceitos, pois o Bandit usa o código um para indicar que encontrou problemas. O adaptador preservou identificador do teste, mensagem, caminho, intervalo, severidade, confiança e CWE informado. A validação rejeitou JSON duplicado, valores não finitos e estruturas incompatíveis, sem alterar a saída original.

\subsection{Adaptador do Semgrep Community Edition}\label{subsec:execucao-semgrep}

O Semgrep Community Edition 1.172.0 utilizou apenas o diretório \texttt{python/} do repositório de regras fixado no \textit{commit} de prefixo \texttt{40b8c63f75dc}. O pacote foi adquirido antes da coleta, teve seu inventário conferido e não consultou o registro remoto durante as análises. O comando efetivo foi:

\begin{verbatim}
semgrep scan --config /opt/regras-semgrep/python --json \
  --metrics=off --jobs 1 \
  --output /saida/bruto.json.part /entrada
\end{verbatim}

O adaptador preservou identificador da regra, mensagem, caminho, intervalo, CWE, severidade e metadados adicionais. Erros parciais e arquivos ignorados foram mantidos em \texttt{avisos.json}, permitindo distinguir execução concluída com avisos de uma falha do processo.

\section{Execução das Condições de IA (C3 e C4)}\label{sec:execucao-ia-desenvolvimento}

\subsection{Adaptadores de sessão}\label{subsec:adaptadores-sessao-ia}

Cursor e Codex foram integrados por adaptadores diferentes, mas receberam o mesmo código, a mesma instrução-base e o mesmo contrato de saída. Cada chamada correspondeu a uma sessão nova e efêmera. O código Python foi serializado em \texttt{payload-codigo.json}, delimitado como dado não confiável e enviado no \textit{prompt} pela entrada padrão. A pasta de trabalho continha somente o esquema e os arquivos mínimos de configuração do cliente; não continha \texttt{oracle/}, identidade original do projeto, respostas anteriores ou resultados de outras condições.

O Cursor foi executado nativamente no Windows pela versão \texttt{2026.09.26-dd393fe}, com o identificador \texttt{gpt-5.6-luna-medium}. O evento de sistema exibiu \texttt{GPT-5.6 Luna 272K Medium}. Como o cliente informava que o isolamento nativo estava disponível apenas em macOS e Linux, a coleta usou \texttt{--sandbox disabled}, \texttt{--mode=ask} e uma política local que permitia \texttt{Read(**)} e negava \texttt{Write(**)} e \texttt{Shell(*)}. O processo foi iniciado na pasta exclusiva da tentativa, e o código seguiu no \textit{prompt}, não como uma árvore navegável no \textit{workspace}. A auditoria dos eventos de C3 e C5, incluindo os pilotos, encontrou zero chamadas de ferramenta; essa constatação reduz, mas não elimina, a limitação do sandbox nativo indisponível no host.

O Codex foi executado pela \texttt{codex-cli 0.154.0}, solicitando \texttt{gpt-5.6-luna}. O comando aplicou sandbox somente leitura e sessão efêmera, ignorou configurações e regras externas e fixou a política de aprovação como \texttt{never}. Os eventos JSONL dessa versão não expuseram o modelo efetivamente servido. Por isso, \texttt{modelo\_exibido} permaneceu nulo e \texttt{modelo\_verificacao} recebeu \path{nao_exposto_jsonl_codex}; o valor solicitado não foi copiado para o campo observado.

O bloco executado foi o principal, com GPT-5.6 Luna nas quatro condições de IA. O bloco complementar previsto como opcional na metodologia, com Cursor Auto e GPT-5.6 Terra no Codex, não integra o conjunto coletado e não foi combinado com a linha de base.

\subsection{Instrução e contrato de saída}\label{subsec:padronizacao-prompts}

O arquivo \texttt{prompt-relatorio-ia-v1.txt} foi compartilhado por C3--C6 e congelado com SHA-256 de prefixo \texttt{1d3e269aa36d}. A instrução solicitou a inspeção do código fornecido, exigiu uma cadeia de exploração plausível e orientou o agente a considerar validações e controles antes de reportar. Também proibiu executar código, testes, \textit{shell}, rede, busca, MCP, Bandit, Semgrep ou qualquer outro analisador, além de proibir consulta ou reconstrução do \textit{ground truth}.

O contrato de saída, congelado com SHA-256 de prefixo \texttt{6e0b2f29d5f3}, exigiu exclusivamente o \textit{array} JSON descrito na Seção~\ref{subsec:esquema-comum-achados}. O controlador validou localmente a sintaxe, as propriedades, os caminhos, as linhas, os hashes, os metadados de modelo e a ausência de eventos proibidos. O argumento \texttt{--output-schema} não foi aplicado ao Codex, pois essa interface exigia um objeto na raiz e criaria tratamento diferente do Cursor. Em ambos os produtos, a saída bruta foi fechada antes da validação comum.

\section{Construção das Condições Híbridas (C5 e C6)}\label{sec:condicao-hibrida-desenvolvimento}

Após C1 e C2, o script \texttt{gerar-alertas-sast.ps1} selecionou a tentativa concluída autoritativa de cada ferramenta e uniu seus achados sem consultar o \textit{ground truth}. A deduplicação considerou caminho, localização, CWE e, quando necessário, identificação da regra e mensagem. Foram observados 1.416 registros antes da deduplicação e 1.410 alertas por repetição depois da remoção de seis duplicatas.

Os 26 arquivos \texttt{ALVO-NNNN-alertas-sast.json} foram congelados em um \textit{lock} que registra SHA-256, quantidade antes e depois da deduplicação e tentativas C1/C2 de origem. Planejamento e execução de C5 e C6 recusam um arquivo que divirja desse \textit{lock}. Dessa forma, Cursor e Codex receberam exatamente os mesmos bytes de alerta para um mesmo alvo.

O \textit{prompt} comum descreveu os alertas como pistas não adjudicadas. O agente deveria confirmá-los com base no código, ignorar falsos positivos e poderia relatar vulnerabilidades sem alerta correspondente. C5 e C6 foram executadas como sessões novas; não reutilizaram contexto nem respostas de C3 e C4. A única diferença experimental planejada entre a condição isolada e sua correspondente híbrida foi, portanto, a presença dos alertas congelados.

\section{Preservação, Auditoria e Preparação das Métricas}\label{sec:adjudicacao-desenvolvimento}

Em C1 e C2, cada tentativa preservou manifestos de estado, os fluxos de saída e erro em formato binário, saída nativa, versões, dados do processo, avisos e saída normalizada. Em C3--C6, foram preservados \textit{prompt} efetivo, \textit{payload}, alertas quando aplicável, fluxo de eventos, \textit{stderr}, resposta bruta, resposta validada e manifesto. O manifesto terminal referencia todos esses arquivos por SHA-256.

A Tabela~\ref{tab:fechamento-operacional-coleta} apresenta o fechamento operacional. ``Tentativas preservadas'' inclui falhas, retentativas e, em C4, tentativas legadas substituídas na seleção. A tabela não representa acurácia e não emprega o \textit{ground truth}.

\begin{table}[htb]
\ABNTEXfontereduzida
\caption{Fechamento operacional das seis condições}
\label{tab:fechamento-operacional-coleta}
\centering
\begin{tabularx}{\textwidth}{|l|X|r|r|r|}
\hline
\textbf{Condição} & \textbf{Ferramenta e repetições} & \textbf{Previstas} & \textbf{Válidas} & \textbf{Tentativas preservadas} \\
\hline
C1--C2 & Bandit e Semgrep, uma execução por alvo & 52 & 52 & 63 \\
\hline
C3 & Cursor, três repetições por alvo & 78 & 78 & 97 \\
\hline
C4 & Codex, três repetições por alvo & 78 & 78 & 90 \\
\hline
C5 & Cursor com alertas, três repetições por alvo & 78 & 78 & 91 \\
\hline
C6 & Codex com alertas, três repetições por alvo & 78 & 78 & 92 \\
\hline
\textbf{Total} & & \textbf{364} & \textbf{364} & \textbf{433} \\
\hline
\end{tabularx}
\legend{Fonte: Elaborado pelo autor a partir das auditorias de fechamento (2026).}
\end{table}

As auditorias foram calculadas a partir dos registros primários e mantiveram separadas as tentativas válidas e inválidas. Duração e tokens foram registrados por tentativa. No Cursor, \texttt{inputTokens}, \texttt{cacheReadTokens} e \texttt{cacheWriteTokens} compuseram os tokens de entrada, enquanto os componentes originais permaneceram nos eventos. No Codex, os contadores expostos pelo JSONL foram preservados sem estimativa. Essas diferenças de observabilidade devem ser consideradas na análise descritiva do Capítulo~\ref{cap:avaliacao}.

O \textit{oracle} somente pode ser introduzido depois do fechamento da coleta. Os achados normalizados e as respostas validadas constituem as entradas da etapa de correspondência com caminho, CWE e tolerância de linha do RealVuln. Não foram usadas correções manuais para completar localização, CWE ou severidade durante a coleta. TP, FP, FN, TN e as métricas derivadas não aparecem nas auditorias operacionais, evitando que uma síntese sem adjudicação seja confundida com resultado experimental.

\section{Verificação do Artefato}\label{sec:verificacao-artefato-desenvolvimento}

Os adaptadores foram exercitados com \textit{fixtures} artificiais externas ao corpus: saída vazia, achado completo, duplicata, JSON inválido, número não finito, erro parcial, CWE ausente e diferentes formatos de caminho e linha. Os testes do controlador também cobriram políticas de aquisição, links e junções, hashes canônicos, preparação, fila, atomicidade, retomada, timeout, códigos de saída e invariantes dos manifestos.

O \textit{gate} de construção reconstruiu a imagem e executou testes unitários, testes do \textit{wrapper}, integrações Docker, verificação das dependências e sondas de isolamento. A sonda observou usuário sem privilégio, ausência de \textit{capabilities}, \texttt{NoNewPrivs=1}, apenas a interface de rede de \textit{loopback}, raiz somente leitura e ausência de corpus e \textit{oracle} na imagem. Uma fumaça descartável sobre um alvo do RealVuln comprovou o fluxo C1/C2 de ponta a ponta e foi arquivada com \texttt{finalidade=fumaca}, fora da coleta principal.

As interfaces de IA possuem testes próprios que simulam eventos e respostas sem realizar chamadas externas. Antes de cada coleta completa, os quatro cenários também passaram por piloto em \texttt{ALVO-0001}. O planejamento em modo seco validou todas as entradas, configurações e hashes sem consumir API. No fechamento, a repetição desse planejamento informou zero chamadas pendentes em C3--C6.

\section{Configuração Efetivamente Observada}\label{sec:registro-configuracao-desenvolvimento}

A Tabela~\ref{tab:registro-configuracao-efetiva} consolida a configuração observada. Os identificadores de \textit{commit} e SHA-256 aparecem abreviados para leitura; os valores completos estão nos \textit{locks}, manifestos e inventários do artefato.

\begin{table}[htb]
\ABNTEXfontereduzida
\caption{Configuração efetiva da implementação e da coleta}
\label{tab:registro-configuracao-efetiva}
\centering
\begin{tabularx}{\textwidth}{|X|X|X|}
\hline
\textbf{Componente} & \textbf{Valor observado} & \textbf{Fonte de evidência} \\
\hline
RealVuln & v1.0; \textit{commit} \texttt{d98e9fc91273}; 26 alvos & \texttt{fontes.lock.json} e resumo do corpus \\
\hline
Corpus rotulado & 817 entradas; 697 vulnerabilidades; 120 armadilhas & Saída do validador oficial e auditoria do \textit{ground truth} \\
\hline
Ambiente SAST & Python 3.12.13; Bandit 1.9.4; Semgrep CE 1.172.0 & \texttt{requirements.lock}, imagem e \texttt{versoes.json} \\
\hline
Regras Semgrep & \textit{commit} \texttt{40b8c63f75dc}; pacote Python verificado & \texttt{fontes.lock.json} e inventário das regras \\
\hline
Cursor & cliente \texttt{2026.09.26-dd393fe}; \texttt{gpt-5.6-luna-medium}; modelo exibido \texttt{GPT-5.6 Luna 272K Medium} & Eventos e manifestos de C3/C5 \\
\hline
Codex & \texttt{codex-cli 0.154.0}; \texttt{gpt-5.6-luna}; modelo servido não exposto no JSONL & Eventos e manifestos de C4/C6 \\
\hline
Instrução de IA & \textit{prompt} \texttt{1d3e269aa36d}; esquema \texttt{6e0b2f29d5f3} & Configurações dos agentes e manifestos \\
\hline
Alertas híbridos & 1.410 alertas por repetição; seis duplicatas removidas & \textit{Lock} de alertas e auditorias C5/C6 \\
\hline
\end{tabularx}
\legend{Fonte: Elaborado pelo autor (2026).}
\end{table}

O host observado na coleta SAST utilizava Git for Windows 2.46.2, WSL 2.0.14, Docker Desktop 4.38.0 e Docker Engine 27.5.1. Como essas versões estavam abaixo das recomendações registradas e incluíam risco conhecido no Docker Desktop, o artefato preserva uma aceitação explícita de risco residual, limitada ao marco SAST e condicionada aos controles de isolamento. A exceção não foi tratada como autorização genérica para outros experimentos.

\section{Desvios e Limitações Observados}\label{sec:desvios-protocolo}

A execução revelou desvios e limitações que não seriam visíveis em uma descrição puramente prospectiva:

\begin{itemize}
\item Nove tarefas C1 falharam inicialmente por incompatibilidade de formato no adaptador do Bandit. A correção foi aplicada em uma imagem identificada separadamente, as tarefas foram reabertas de forma explícita e concluídas na tentativa 3. Saídas antigas, transições da fila e ambos os identificadores de imagem foram preservados.

\item C4 foi a primeira condição de IA concluída. Setenta unidades selecionadas provinham de uma versão anterior do executor, sem o campo \texttt{modelo\_verificacao} e sem a validação atual de caminho e linha. Oito unidades foram reprocessadas; seis manifestos legados antes marcados como concluídos foram substituídos na seleção, mas não apagados nem reescritos. A auditoria registra 90 tentativas e os dois hashes de configuração históricos.

\item A validação estrita rejeitou respostas de IA que mencionavam arquivo inexistente ou linha fora do arquivo. Foram preservadas 19 falhas de formato em C3, seis em C4, 13 em C5 e 14 em C6. Novas sessões produziram as tentativas seguintes até completar as 78 unidades válidas de cada condição. Uma falha de formato nunca foi convertida em lista vazia ou ausência de achados.

\item No Cursor nativo para Windows, o sandbox de sistema permaneceu desativado. A mitigação combinou área exclusiva, negação de escrita e \textit{shell}, entrada de código por \textit{stdin}, rejeição de eventos proibidos e auditoria posterior. A permissão \texttt{Read(**)} representa capacidade relativa ao \textit{workspace}, não prova de leitura; embora nenhum \texttt{tool\_call} tenha sido observado, essa diferença em relação ao sandbox somente leitura do Codex permanece uma ameaça à validade.

\item A CLI do Codex não expôs no JSONL o modelo efetivamente servido. A implementação preservou essa ausência em vez de inferir que o modelo solicitado foi necessariamente o observado. Em parte do histórico de C4, nem mesmo o campo explicativo atual estava disponível; essa limitação foi mantida nos manifestos legados.

\item Dez alvos e o pacote local de regras Semgrep não puderam ser redistribuídos integralmente. Essa restrição não alterou as entradas da coleta local; afetou apenas a forma de publicação. URLs, revisões, inventários e scripts de reconstrução foram mantidos para permitir auditoria.
\end{itemize}

Ao final, todas as unidades previstas do bloco principal possuíam uma tentativa selecionada e válida, sem tarefas pendentes. A preservação dos desvios, em vez da sua remoção retrospectiva, mantém a rastreabilidade entre configuração, chamada externa, resposta bruta, validação e conjunto escolhido para a avaliação.
